\documentclass[11pt]{article}
\usepackage[a4paper,margin=1in]{geometry}
\usepackage{amsmath,amssymb,amsthm,booktabs}
\newtheorem{theorem}{Theorem}
\newtheorem{proposition}[theorem]{Proposition}
\newtheorem{lemma}[theorem]{Lemma}
\newcommand{\Q}{\mathbf Q}
\title{A streamlined computer-assisted proof of\\
Sharipov's second cuboid conjecture}
\author{Computational research draft}
\date{August 11, 2026}

\begin{document}
\maketitle

\begin{abstract}
We prove that Sharipov's second cuboid polynomial is irreducible for every
pair of coprime positive integers $p\ne q$.  The proof reduces factorization
to the $1+4$, $2+3$, and symmetric $5+5$ cases, and compresses the completed
$2+3$ argument into one lossless reduction to
\[
 C_3:\quad r^2=4u^8-8u^2+5
\]
and a direct classification of $C_3(\Q)$.  The existing exact two-descent
supplies the quotient rank.  The unique $3$-adic Chabauty differential has one extra
$\Q_3$-zero in each of four affine disks; an integral Mordell--Weil lattice
argument excludes all four.  Exact certificates retain every denominator
and exceptional fibre.  No bounded search, finite sieve as a completeness
premise, or open direct-normalization claim is used.
\end{abstract}

\section{Factorization reduction}

For $s=(p/q)^2$, put
\[
\begin{aligned}
P_s(x)={}&x^5+(2+s)(3-2s)x^4
 +(1+10s+4s^2-14s^3+s^4)x^3\\
&-s(1-14s+4s^2+10s^3+s^4)x^2\\
&-s^3(1+2s)(-2+3s)x-s^5
\end{aligned}
\]
and
\[
 Q_{p,q}(t)=q^{20}P_s(t^2/q^4).
\]
The admissible domain is $p,q>0$, $\gcd(p,q)=1$, and $p\ne q$; hence
$s>0$ and $s\ne1$.

\begin{proposition}[Complete factorization reduction]\label{prop:factor}
It is enough to exclude a rational linear factor of $P_s$, a rational
quadratic factor of $P_s$, and a symmetric factorization
\[
 P_s(u^2)=R(u)(-R(-u))
\]
with $R$ irreducible of degree five.
\end{proposition}

\begin{proof}
The scaling $u=t/q^2$ and the nonzero factor $q^{20}$ preserve
reducibility.  If the monic quintic $P_s$ is reducible, one of its
irreducible factors has degree at most two, giving the first two cases.
If $P_s$ is irreducible and $\alpha$ is a root, Capell's composition
criterion says that $P_s(u^2)$ is reducible precisely when $\alpha$ is a
square in $\Q(\alpha)$.  Writing $\alpha=\beta^2$, one has
$\Q(\beta)=\Q(\alpha)$, so the minimal polynomial $R$ of $\beta$ has degree
five and conjugating $\beta$ to $-\beta$ gives the displayed factorization.
Conversely that factorization is reducible.  Finally, monicity and Gauss's
lemma identify rational and integral factorizations of $Q_{p,q}$.
\end{proof}

\section{The linear--quartic case}

\begin{theorem}[Linear-factor exclusion]\label{thm:linear}
For positive rational $s\ne1$, the quintic $P_s$ has no rational root.
\end{theorem}

\begin{proof}
The exact reciprocal quotient sends a rational root to a rational point on
\[
 C_1:\quad w^2=t^5+21t^4+26t^3+10t^2+5t+1.
\]
A Schaefer--Stoll descent in
$\Q\times\Q[z]/(z^4-2z-1)$, with certified class and unit groups, proves
\[
 \dim_{\mathbf F_2}\operatorname{Sel}_2(J_1/\Q)=2,\qquad
 \operatorname{rank}J_1(\Q)=1,\qquad
 \operatorname{Sha}(J_1/\Q)[2]=0.
\]
At $5$, the six known rational points reduce bijectively to
$C_1(\mathbf F_5)$.  An annihilating differential reduces to
\[
 (1+t)\frac{dt}{2w}.
\]
It is nonzero in five disks.  In the remaining disk at $(-1,0)$, the
parameter $w=5z$ gives the exact Strassmann form
\[
 \frac{G(5z)}{25}
 =z\bigl(\text{\rm unit}+5z^2\mathbf Z_5[[z^2]]\bigr).
\]
Thus
\[
 C_1(\Q)=\{\infty,(-1,0),(0,\pm1),(1,\pm8)\}.
\]
The inverse quotient sends these points only to $s=1$, $s=-1$, or poles.
The remaining denominator $t^2+6t+1$ has discriminant $32$ and no rational
zero.  Hence no admissible $s$ occurs.  The descent, local series, quotient
identities, and lift table are replayed by the certificates in
\texttt{research/linear\_factor\_completion/}.
\end{proof}

\section{A lossless reduction of the quadratic--cubic case}

\begin{lemma}[Signed-square constant]\label{lem:signed}
If $X^2+AX+B$ is an integral monic factor of the quintic in $X=t^2$, then
\[
 B=\sigma d^2,\qquad \sigma\in\{+1,-1\},\quad d\in\mathbf Z.
\]
\end{lemma}

\begin{proof}
For a prime $\ell\mid pq$, put $k=v_\ell(pq)$.  Direct coefficient
valuations give the lower Newton polygon with vertices
\[
 (0,10k),(2,2k),(3,0),(5,0),
\]
including $\ell=2,3$.  The root valuations are
$4k,4k,2k,0,0$, so the product of any two roots has even valuation.
Since $B$ divides $-p^{10}q^{10}$, all valuations of $B$ are even.  The
single sign $\sigma$ is physical and global.
\end{proof}

\begin{theorem}[Consolidated $2+3$ reduction]\label{thm:reduction}
Every admissible rational quadratic factor of $P_s$ determines a rational
solution of
\[
 C_3:\quad r^2=4u^8-8u^2+5,                         \tag{1}
\]
and hence a point
\[
 (a,R)=(u^2,ur)
\quad\hbox{on}\quad
 D_+:\ R^2=4a^5-8a^2+5a.                            \tag{2}
\]
Every source, base-locus, denominator-zero, and exceptional fibre omitted
by the generic formulas is explicitly finite and inadmissible.
\end{theorem}

\begin{proof}
We give the exact chain while suppressing large coordinates from the
statement.  Write $s=z^2$, normalize the constant of the quadratic factor as
$b=\sigma z^4k^2$, and write its linear coefficient as $z^2ky$.
Substitution in the two exact remainder equations and elimination of $y$
gives an irreducible $71$-term curve
\[
 Q_\sigma(z,k)=0.
\]
The penultimate subresultant is $L_1y+L_0$ and recovers
$y=-L_0/L_1$.  The resultant of $Q_\sigma$ with the non-boundary factor of
$L_1$ is
\[
\begin{aligned}
2^{20}z^{24}(z-1)^{12}(z+1)^{12}
 &(z^2-z+1)^4(z^2+z+1)^4(z^2+1)^{12}\\
 &\cdot F_{32}(z)^4,
\end{aligned}
\]
where $F_{32}$ is the explicit degree-$32$ factor stored in the compact-model
certificate; it is monic with constant $1$ and $F_{32}(\pm1)=512$.  The
quadratic factors have negative discriminant.  Thus the only rational
exceptional values are $z=0,\pm1$, namely $s=0,1$.  This proves that the
compact model and inverse are lossless on the admissible locus.

After putting $S=z^2$, $K=k^2$ and applying two Cayley transforms, the
simultaneous reciprocal involution has fixed generators
$\mathcal U,\mathcal V$.  Both physical signs give the same irreducible
$29$-term curve $C_8(\mathcal U,\mathcal V)=0$ and require
\[
 \mathcal U\mathcal V,\qquad 1-\mathcal U\mathcal V       \tag{3}
\]
to be nonzero rational squares.  Exact quartic functions give maps in both
directions between the admissible normalization of $C_8$ and
\[
 E_1:\quad y^2=x^3-5x+4.                                \tag{4}
\]
Substitution, linear subresultants, and multiplication-back verify both
round trips.  The factored source and inverse denominator resultants leave
only the fibres in Table~\ref{tab:exceptions}; each is a boundary,
nonrational, sign-incompatible, or already fails (3).  In particular the
negative source chart's potentially lost radial fibre has discriminant
$2^8\cdot5$, and the removable elliptic base point has odd $2$-adic
valuation of $\mathcal U\mathcal V$.  No generic open-set assertion replaces
this audit.

Translate (4) by its rational two-torsion point, obtaining coordinates
$(X,Y)$, and put
\[
L=2X+Y-2,\quad M=X^2+2Y-1,\quad
\kappa=\frac{Y+2}{X},\quad h=X+1.
\]
Exact reduction in the function field gives
\[
 XM=(Y+2)L,\qquad
 1-\mathcal U\mathcal V
 =2(X+1)LMg^2                                      \tag{5}
\]
for a recorded rational function $g$.  The factor $2$ in (5) is essential.
Away from the separately audited zeros and poles, (3) implies
\[
 \xi^2=2\kappa h,\qquad
 h^2-(\kappa^2+2)h+\kappa^2+4\kappa-4=0.           \tag{6}
\]

For every odd prime, the Newton polygon of the second equation and the first
equation in (6) show that $v(\kappa)$ is even.  At $2$, the discriminant
$\kappa^4-16\kappa+20$ excludes $v_2(\kappa)=0$ and
$v_2(\kappa)\ge2$; negative valuation has root valuations $0$ and
$2v_2(\kappa)$.  Hence
\[
 \kappa=2\epsilon u^2,\qquad h=\epsilon v^2,\qquad
 \epsilon\in\{+1,-1\}.                               \tag{7}
\]
This valuation sign is unrelated to the physical sign $\sigma$.
For $\epsilon=-1$, clearing denominators gives
\[
 W^2=-(v^2+2)(v^4+v^2-4).
\]
The coprime parity cases give $3$, $6$, or $7$ modulo $8$, or an odd
$2$-adic valuation, so there is no rational point.  For $\epsilon=1$, the
discriminant of (6) is exactly (1), and (2) follows.

The zeros and poles divided out above are precisely those in
Table~\ref{tab:exceptions}.  Direct evaluation uses the original characters,
so even the fibre $L=M=0$ is not lost.
\end{proof}

\section{Direct $3$-adic Chabauty on $C_3$}

Let $J_+=\operatorname{Jac}(D_+)$,
\[
 D=[(1,1)-\infty],\qquad T=[(0,0)-\infty]\in J_+(\Q).
\]

\begin{lemma}[Rank and annihilator]\label{lem:rank}
One has
\[
 \operatorname{rank}J_+(\Q)=1,\quad
 J_+(\Q)_{\rm tors}=\langle T\rangle\simeq\mathbf Z/2\mathbf Z,
\]
and $D,T$ are independent modulo $2$.  Moreover
\(\operatorname{rank}\operatorname{Jac}(C_3)(\Q)=2\), and its
$3$-adic Chabauty annihilator is the line generated by
\[
 \eta=(\lambda-u^2)\frac{du}{r},\qquad
 \lambda\equiv7\pmod {27},\quad
 v_3(\lambda-1)=1,\quad v_3(\log_0D)=2.             \tag{8}
\]
\end{lemma}

\begin{proof}
For the genus-two rank input, a complete Schaefer--Stoll descent in
\[
 \Q\times\Q[z]/(z^4-16z+20),\qquad S=\{2,5,17\},
\]
has exact survivor chain
\[
 512\longrightarrow64\longrightarrow8\longrightarrow4.
\]
The final classes are \(0,D,T,D+T\), and the condition at \(17\) is active.
Together with exact good-reduction arithmetic, this proves
\[
 \operatorname{rank}J_+(\Q)=1,\qquad
 J_+(\Q)_{\rm tors}=\langle T\rangle,\qquad
 D,T\text{ independent modulo }2.
\]
The full class/unit, norm, and local-image certificate is replayed exactly;
no primitive generator is assumed.

The two involution quotients of $C_3$ give
\[
 \operatorname{Jac}(C_3)\sim_\Q E_1\times J_+,\qquad
 H^0(C_3,\Omega^1)=
 \left\langle\frac{u\,du}{r}\right\rangle\oplus
 \left\langle\frac{du}{r},\frac{u^2du}{r}\right\rangle.
\]
The elementary $2$-isogeny descent gives
\(\operatorname{rank}E_1(\Q)=1\).  At \(3\), the quotient isogeny has
prime-to-\(3\) degree, both factors have good reduction, and the logarithm
kernel is torsion.  The two non-torsion factor logarithms are therefore
nonzero in the displayed independent tangent summands.  Thus the annihilator
is one-dimensional and comes from $J_+$.  Exact Cantor arithmetic for $11D$,
followed by
divided formal integrals of $da/R$ and $a\,da/R$, gives
\[
 \lambda=\frac{\log_1D}{\log_0D}\equiv7\pmod {27},
 \qquad v_3(\log_0D)=2.
\]
The support discriminant has valuation $4$, and every omitted primitive
degree $n\ge4$ has valuation at least
$n-1-v_3(n)\ge3$.  This proves (8) without assuming that $D$ is primitive.
\end{proof}

\begin{theorem}[Direct rational-point theorem]\label{thm:c3}
\[
 C_3(\Q)=
 \{\infty_+,\infty_-,(1,1),(1,-1),(-1,1),(-1,-1)\}. \tag{9}
\]
\end{theorem}

\begin{proof}
The defining polynomial is squarefree modulo \(3\), so \(C_3\) has good
reduction there.  Exact enumeration gives $\#C_3(\mathbf F_3)=6$, and the
points in (9) occupy the six residue disks.  At either infinity put
\(t=1/u\), \(\rho=rt^4\).  Then
\[
 \rho^2=4-8t^6+5t^8,\qquad
 \eta=(1-\lambda t^2)\frac{dt}{\rho}.
\]
After $t=3z$, its primitive divided by $3$ is $z$ times a unit; the
uniform tail bound is $n-1-v_3(n)\ge1$ for $n\ge2$.  Hence each infinity
disk contains only its known point.

In an affine disk centered at $(u,r)=(s,r_0)$, put $u=s+3z$.  The exact
primitive $F$ satisfies, up to a unit,
\[
 \frac{F(3z)}9\equiv2z-sz^2\pmod3.                 \tag{10}
\]
The coefficient bounds are
\[
 n-1-v_3(n)\ge1\ (n\ge2),\qquad
 n-2-v_3(n)\ge1\ (n\ge4).
\]
Thus every affine disk has two simple $\Q_3$ Coleman zeros: the known
$z=0$ and one extra zero.  Plain Chabauty on the full disks is therefore
insufficient.

To remove them, after a torsion shift write
\[
 J_+(\Q)=\mathbf ZG\oplus\langle T\rangle,\quad D=mG,\quad
 \#J_+(\mathbf F_7)=42,\quad\operatorname{ord}(\bar D)=21.
\]
Modulo-$2$ independence makes $m$ odd.  Since
\(\operatorname{ord}(\bar G)/\gcd(\operatorname{ord}(\bar G),m)=21\),
the alternative \(\operatorname{ord}(\bar G)=42\) would make \(m\) even.
Thus \(\gcd(m,21)=1\) and \(v_3(\log_0G)=2\).

For a rational point $Q$ in an affine disk, map to $D_+$ by
\(\pi(u,r)=(a,R)=(u^2,ur)\), and let $P$ be the known center.  If
$x=a(\pi(Q))-1$, then for some \(n\in\mathbf Z\) and \(e\in\{0,1\}\),
\[
\begin{gathered}
 \relax[\pi(Q)-\pi(P)]=nG+eT,\qquad \log T=0,\\
 \log_0[\pi(Q)-\pi(P)]=n\log_0G\in9\mathbf Z_3,\qquad
 \int_0^x\frac{da}{R}=x\cdot\mathbf Z_3^\times.
\end{gathered}
\]
Here \(da/du=2u\) and \(da/R\) are units in these disks.
Thus $x\in9\mathbf Z_3$.  On writing $x=9z$, the annihilating primitive
divided by $27$ reduces to a nonzero scalar times $z$ modulo $3$; the tail
bounds are \(2n-2-v_3(n)\ge1\) and \(2n-3-v_3(n)\ge1\) for \(n\ge2\).
Strassmann gives $z=0$, hence $a=1$, $u=\pm1$, and $r=\pm1$.  This treats
all four affine disks and proves (9).
\end{proof}

\section{Exceptional lifts and conclusion}

\begin{table}[ht]
\centering
\small
\begin{tabular}{p{.20\textwidth}p{.31\textwidth}p{.39\textwidth}}
\toprule
Stage & Rational candidate & Exact disposition\\
\midrule
Compact inverse & $z=0,\pm1$ & $s=0$ or $s=1$.\\
Reciprocal charts & $S,K=0,\pm1$ and the negative radial zero &
Boundary; a constant-one quartic has no rational root; the radial
quadratic has discriminant $2^8\cdot5$.\\
$C_8\leftrightarrow E_1$ & $\mathcal U=0,\pm1,1/2$; $x=2$, $x^2=5$,
$\pm4P$ for $P=(0,2)$, infinity & No rational inverse, boundary or wrong sign,
irreducible fibre, or a common character with odd $2$-adic valuation;
infinity has $\mathcal U\mathcal V=1/2$.\\
Character identity & $X=0,1,3,-1$ and zeros/poles of $g,L,M$ &
Direct values are boundary or nonsquare.  At $(X,Y)=(1,0)$,
$L=M=0$ but the undivided character is $1/2$.\\
Final $C_3$ points & $\infty_\pm$ and $(u,r)=(\pm1,\pm1)$ &
The infinities are boundary.  The affine points map to
$(a,R)=(1,\pm1)$; their translated elliptic fibres have undivided
common-character value $1/2$ or $0$.  The omitted $a=0$ lift would require
$r^2=5$.\\
\bottomrule
\end{tabular}
\caption{Exact exceptional-fibre audit for Theorem~\ref{thm:reduction}.}
\label{tab:exceptions}
\end{table}

\begin{theorem}[Sharipov's second cuboid conjecture]
For coprime positive integers $p\ne q$, the polynomial $Q_{p,q}(t)$ is
irreducible over $\Q$.
\end{theorem}

\begin{proof}
By Theorem~\ref{thm:c3}, a point arising from
Theorem~\ref{thm:reduction} is one of the six points in (9).
Table~\ref{tab:exceptions} excludes every lift: in particular, the four
affine points map to $a=1$, where the two translated elliptic points are
$(1,0)$ and $(3,4)$, with undivided common-character values $1/2$ and $0$.
Thus no $2+3$ factor exists.

Theorem~\ref{thm:linear} excludes $1+4$.  For the symmetric $5+5$ case, the
pinned source archive of arXiv:2601.04240 has SHA--256
\[
\texttt{96d7e4b7d3cb75af238a0ca8cd143a8d555daf6ac9e8c2cc94a099b5aea04760}.
\]
The replay extracts its complete second code block and exactly recomputes the
elimination polynomial, discriminant factorization, rational-root tests,
positive-root isolations, and all seven Sturm counts.  Every admissible
positive parameter interval has zero real candidate roots; the special
parameter $s=1$ is factored separately and is inadmissible.  Hence the
symmetric $5+5$ split is impossible.  Proposition~\ref{prop:factor} now
excludes every proper factorization.
\end{proof}

\begingroup\small
\paragraph{Reproducibility and dependency change.}
The proof-critical commands are:
\begin{verbatim}
PYTHON_BIN=.venv/bin/python bash scripts/reproduce.sh
PYTHON_BIN=.venv/bin/python bash scripts/replay_5plus5.sh
tectonic paper/clean2.tex
\end{verbatim}
The \(512\to64\to8\to4\) rank descent remains; direct \(p=3\) Chabauty
replaces the \(p=23\) sieve, \(p=7\) square tables, \(42D\) logarithm, Coleman
differential, and coefficient comparisons.  A valid Richelot-isogeny descent
is replayed as an alternative, but is not a dependency because its complete
certificate is larger.  Reductions at \(3,7\) remain for rank, torsion, and
\(3\)-saturation.  C025, finite sieves, fingerprints, genus-$33$ covers, and
bounded searches are not dependencies.
\endgroup

\end{document}
